\documentclass[11pt,a4paper]{article}
\usepackage[utf8]{inputenc}
\usepackage[T1]{fontenc}
\usepackage[french]{babel}
\usepackage{amsmath}
\usepackage{amssymb}
\usepackage{graphicx}
\usepackage{hyperref}
\usepackage{listings}
\usepackage{xcolor}
\usepackage{geometry}
\usepackage{booktabs}
\usepackage{enumitem}
\geometry{margin=1in}

% Style pour le code Python
\lstset{
    language=Python,
    basicstyle=\ttfamily\small,
    keywordstyle=\color{blue},
    commentstyle=\color{gray},
    stringstyle=\color{red},
    showstringspaces=false,
    breaklines=true,
    frame=single,
    numbers=left,
    numberstyle=\tiny\color{gray}
}

\title{Réponses aux questions de présentation\\Pipeline d'apprentissage par renforcement}
\author{}
\date{8 février 2026}

\begin{document}
\maketitle
\tableofcontents
\newpage

\section{Question 1 — Pourquoi l'Average Loss est-elle naturellement négative ?}

\subsection{Explication}

L'Average Loss (avg\_loss) est une \textbf{somme pondérée} de plusieurs composantes :

\begin{equation}
L = 1{,}0 \cdot L_{\text{policy}} + 0{,}3 \cdot L_{\text{contrastive}} + 0{,}5 \cdot L_{\text{value}} - 0{,}1 \cdot H + \ldots
\end{equation}

La composante dominante est la \textbf{policy loss}, définie comme :

\begin{equation}
L_{\text{policy}} = -\frac{1}{B}\sum_{i} \log(\cos(\hat{a}_i, a_i)) \cdot \hat{A}_i
\end{equation}

où $\hat{A}_i = \frac{r_i - \bar{r}}{\sigma_r + 10^{-8}}$ est l'avantage normalisé.

\subsection{Mécanisme}

Quand l'agent s'améliore :
\begin{itemize}
    \item La cosine similarity $\cos(\hat{a}, a)$ augmente (se rapproche de 1)
    \item $\log(\cos)$ se rapproche de 0 (depuis des valeurs négatives)
    \item Multiplié par un avantage positif $\hat{A} > 0$, on obtient une valeur négative
    \item Le signe négatif devant rend la policy loss \textbf{de plus en plus négative}
\end{itemize}

Dans le baseline :
\begin{itemize}
    \item Policy loss : $\approx -6{,}0$
    \item Contrastive loss : $\approx 3{,}3$ (faible)
    \item Value loss : $\approx 0{,}003$ (très faible)
\end{itemize}

Résultat :
\begin{equation}
L \approx 1{,}0 \times (-6{,}0) + 0{,}3 \times 3{,}3 + 0{,}5 \times 0{,}003 \approx -5{,}0
\end{equation}

\subsection{Est-ce normal en RL ?}

\textbf{Oui, c'est normal.} Contrairement au supervised learning où la loss est positive et descend vers 0, en RL la loss est un \textbf{objectif à maximiser} (la récompense attendue). On minimise $-J(\theta)$, ce qui revient à minimiser une quantité qui devient de plus en plus négative.

\textbf{Une loss qui descend (plus négative) = un agent qui s'améliore.}

\section{Question 2 — Rôle de la contrastive loss}

\subsection{Définition et rôle}

La \textbf{contrastive loss (InfoNCE)} apprend à l'agent à :
\begin{itemize}
    \item Produire des embeddings qui \textbf{ressemblent au bon prompt} (le positif)
    \item Être \textbf{éloigné des mauvais prompts} (les négatifs)
    \item Structurer l'espace de représentation
\end{itemize}

\subsection{Formule}

\begin{equation}
L_{\text{InfoNCE}} = -\log \frac{e^{\text{sim}(p, a^+)/\tau}}{e^{\text{sim}(p, a^+)/\tau} + \sum_j e^{\text{sim}(p, a^-_j)/\tau}}
\end{equation}

où :
\begin{itemize}
    \item $p$ : prédiction de l'agent
    \item $a^+$ : prompt positif (bonne récompense)
    \item $a^-_j$ : prompts négatifs (mauvaises récompenses)
    \item $\tau$ : température (contrôle la netteté de la distribution)
\end{itemize}

\subsection{Quand diminue-t-elle ?}

\begin{itemize}
    \item L'agent produit des prédictions très proches du bon prompt (sim élevée avec le positif)
    \item Les prédictions sont éloignées des mauvais prompts (sim faible avec les négatifs)
    \item Les embeddings sont bien \textbf{séparés} dans l'espace latent
    \item Le ratio $\frac{\text{positif}}{\text{positif + négatifs}}$ augmente
\end{itemize}

\subsection{Quand augmente-t-elle ?}

\begin{itemize}
    \item Les embeddings des candidats sont tous proches les uns des autres
    \item L'agent ne distingue pas le bon prompt du mauvais
    \item Le nombre de négatifs augmente (plus de distracteurs dans le dénominateur)
    \item La température $\tau$ est trop basse (amplifie les petites différences)
    \item Les embeddings ne sont pas normalisés correctement
\end{itemize}

\section{Question 3 — Rôle de la policy loss}

\subsection{Définition et rôle}

La \textbf{policy loss} est le cœur de l'apprentissage RL. Elle ajuste les paramètres du réseau pour que l'agent :
\begin{itemize}
    \item \textbf{Produise plus souvent} les actions qui ont mené à de bonnes récompenses
    \item \textbf{Évite} celles qui ont mené à de mauvaises récompenses
\end{itemize}

\subsection{Formule (REINFORCE avec baseline)}

\begin{equation}
L_{\text{policy}} = -\log(\cos(\hat{a}, a)) \cdot \hat{A}
\end{equation}

où $\hat{A} = \frac{r - \bar{r}}{\sigma_r}$ est l'avantage normalisé.

\subsection{Quand devient-elle plus négative (= amélioration) ?}

\begin{itemize}
    \item La cosine similarity entre la prédiction et l'action choisie augmente
    \item $\log(\cos)$ se rapproche de 0
    \item L'avantage est positif (récompense supérieure à la moyenne)
    \item Combiné : $-(0^-) \times (+) = $ valeur négative qui se creuse
\end{itemize}

\subsection{Quand augmente-t-elle (= problème) ?}

\begin{itemize}
    \item L'agent prédit des vecteurs éloignés des bonnes actions
    \item $\log(\cos)$ est très négatif (cosine similarity faible)
    \item L'avantage est bruité ou mal estimé
    \item Le critique (value function) est imprécis
    \item Les mises à jour vont dans la mauvaise direction
\end{itemize}

\section{Question 4 — Comment détecter l'overfitting en RL ?}

\subsection{Contexte}

En RL, on ne parle pas d'overfitting au sens classique (train vs test), mais d'un phénomène similaire : \textbf{la politique se spécialise trop sur les patterns rencontrés} et perd en généralisation.

\subsection{Observation à analyser}

« Légère régression de la reward en fin de training (premiers 50 : 0{,}725 $\rightarrow$ derniers 50 : 0{,}651, -10{,}2\%). »

\subsection{Hypothèse 1 — Overfitting}

L'agent a mémorisé des associations « ce type de phrase $\rightarrow$ ce type de prompt » qui fonctionnaient au début mais ne se généralisent pas aux nouvelles phrases.

\textbf{Indices pour confirmer :}
\begin{itemize}
    \item L'entropie de la politique baisse (exploration réduite)
    \item L'epsilon très bas en fin de training (0{,}05)
    \item Les mêmes prompts sont sélectionnés de plus en plus souvent
    \item La value loss augmente (le critique ne suit plus)
\end{itemize}

\subsection{Hypothèse 2 — Changement de distribution}

Les phrases 501--600 sont intrinsèquement plus difficiles :
\begin{itemize}
    \item Plus courtes (moins de contexte acoustique)
    \item Vocabulaire rare
    \item Structures syntaxiques inhabituelles
\end{itemize}

Les statistiques par bloc le confirment partiellement :
\begin{itemize}
    \item Bloc 201--300 : CER = 0{,}344 (pire)
    \item Bloc 401--500 : CER = 0{,}290 (meilleur)
    \item Bloc 501--600 : CER = 0{,}299 (légère remontée)
\end{itemize}

\subsection{Comment trancher ?}

Tester le modèle sauvegardé à l'itération 400 sur les phrases 501--600 :
\begin{itemize}
    \item Si le CER reste meilleur $\rightarrow$ c'est de l'overfitting
    \item Si le CER est similairement mauvais $\rightarrow$ changement de distribution
\end{itemize}

\section{Question 5 — Embeddings « mal séparés », positifs et négatifs}

\subsection{Positif et négatif dans le contrastive learning}

\begin{itemize}
    \item \textbf{Candidat positif ($a^+$)} : le prompt qui a donné la \textbf{meilleure récompense} (bon CER). C'est celui que l'agent devrait apprendre à sélectionner.
    \item \textbf{Candidats négatifs ($a^-$)} : tous les autres prompts du batch avec des récompenses \textbf{inférieures} à la baseline. Ce sont ceux dont l'agent devrait s'éloigner.
\end{itemize}

\subsection{« Embeddings mal séparés » — Explication visuelle}

Imagine un espace 256D. 

\textbf{Bien séparés :} le bon prompt (positif) forme un cluster éloigné des mauvais prompts (négatifs). L'agent peut facilement les distinguer — c'est comme choisir la balle rouge parmi des balles bleues.

\textbf{Mal séparés :} tous les candidats sont regroupés au même endroit dans l'espace. Le positif et les négatifs ont des cosine similarities presque identiques vis-à-vis de la prédiction. C'est comme chercher la balle rouge parmi des balles orange — très difficile.

\subsection{« Le nombre de candidats négatifs est plus élevé »}

Dans le dénominateur de l'InfoNCE :
\begin{equation}
\sum_j e^{\text{sim}(p, a^-_j)/\tau}
\end{equation}

Si on passe de 5 négatifs à 63, le dénominateur grossit énormément, ce qui rend le ratio $\frac{\text{positif}}{\text{positif} + \text{négatifs}}$ beaucoup plus petit.

Résultat : $-\log(\text{petit nombre})$ = grande valeur de loss.

\subsection{« Leur similarité est trop proche »}

Exemple numérique : positif sim = 0{,}80, négatifs sim = 0{,}78, 0{,}79, $\ldots$

\begin{equation}
p(\text{positif}) = \frac{e^{0{,}8/0{,}1}}{e^{0{,}8/0{,}1} + e^{0{,}78/0{,}1} + e^{0{,}79/0{,}1} + \ldots}
\end{equation}

Si les exponentielles sont presque égales avec 63 négatifs :
\begin{equation}
p \approx \frac{e^8}{e^8 + 63 \times e^{7{,}85}} \approx \frac{1}{1 + 63 \times e^{-0{,}15}} \approx \frac{1}{55} \approx 0{,}018
\end{equation}

La loss : $-\log(0{,}018) \approx 4{,}0$ (proche du cas aléatoire $-\log(1/64) = 4{,}16$).

\subsection{« La contrastive loss a explosé d'un facteur ×58 »}

\begin{table}[h]
\centering
\begin{tabular}{lcc}
\toprule
\textbf{Expérience} & \textbf{Contrastive Loss} & \textbf{Signification} \\
\midrule
Baseline (260326) & 3{,}3 & Embeddings bien séparés \\
Pipeline v2 (261011) & 192{,}1 & Embeddings quasi-indistinguables \\
\textbf{Facteur} & \textbf{×58} & \\
\bottomrule
\end{tabular}
\end{table}

Le modèle v2 voit tous les candidats comme interchangeables. Le signal d'apprentissage de la contrastive loss est quasi-nul.

\section{Question 6 — Température et normalisation}

\subsection{La température $\tau$}

Dans la formule $e^{\text{sim}/\tau}$, la température contrôle la \textbf{netteté} de la distribution :

\begin{itemize}
    \item \textbf{$\tau$ petit (ex: 0{,}05)} : distribution très piquée. Même une petite différence de similarité (0{,}80 vs 0{,}82) produit un écart énorme dans les probabilités. Le modèle est très discriminant mais les gradients peuvent être instables.
    \item \textbf{$\tau$ grand (ex: 0{,}5)} : distribution plate. Les similarités sont « lissées », le modèle traite les candidats de manière plus uniforme. Plus stable mais apprend plus lentement.
\end{itemize}

\subsection{Impact du changement de température}

Si la température a changé entre le baseline ($\tau = 0{,}1$ fixe) et le v2 (température adaptative entre 0{,}05 et 0{,}5), cela modifie directement l'échelle de la contrastive loss.

\textbf{Exemple numérique :}
\begin{align}
\text{Avec } \tau = 0{,}1 : \quad & e^{0{,}8/0{,}1} = e^{8} \approx 2981 \\
\text{Avec } \tau = 0{,}5 : \quad & e^{0{,}8/0{,}5} = e^{1{,}6} \approx 5
\end{align}

Le ratio change dramatiquement, affectant la loss finale.

\subsection{La normalisation}

Les embeddings sont L2-normalisés ($\|e\| = 1$), ce qui contraint la cosine similarity dans $[-1, 1]$.

\textbf{Si la normalisation change :}
\begin{itemize}
    \item Pas de normalisation $\rightarrow$ les vecteurs ont des normes arbitraires
    \item $\text{sim}/\tau$ peut prendre des valeurs beaucoup plus grandes ou plus petites
    \item La loss explose car les exponentielles sont décalées
\end{itemize}

\section{Question 7 — Pourquoi la contrastive loss varie d'une phrase à l'autre}

\subsection{Raisons fondamentales}

\begin{enumerate}
    \item \textbf{Les candidats changent} : Pour chaque phrase cible, FAISS récupère un ensemble différent de prompts voisins. Certaines phrases ont des prompts très distincts (facile $\rightarrow$ loss faible), d'autres ont des prompts tous similaires (difficile $\rightarrow$ loss élevée).
    
    \item \textbf{Le positif change} : La qualité du meilleur prompt varie. Si le meilleur prompt donne une similarité de 0{,}9 avec la prédiction, la loss sera basse. Si c'est 0{,}3, elle sera haute.
    
    \item \textbf{La composition du replay buffer} : Quand on entraîne sur un batch, les récompenses des échantillons déterminent qui est positif/négatif. Un batch avec des récompenses très homogènes (tous autour de 0{,}6) rend la séparation plus floue qu'un batch avec des récompenses contrastées (0{,}1 et 0{,}9).
\end{enumerate}

\subsection{Analogie intuitive}

C'est comme un examen QCM :
\begin{itemize}
    \item \textbf{Questions faciles} : réponse évidente, distracteurs éloignés $\rightarrow$ loss basse
    \item \textbf{Questions difficiles} : options très proches les unes des autres $\rightarrow$ loss haute
\end{itemize}

La difficulté change \textbf{à chaque question} (= à chaque phrase).

\section{Question 8 — La value loss qui augmente vs le plot qui descend}

\subsection{Distinction importante}

Il faut bien distinguer \textbf{deux choses différentes} :

\begin{itemize}
    \item \textbf{Le plot de loss} : montre l'Average Loss (somme pondérée de toutes les composantes). Cette somme descend car la policy loss et la contrastive loss descendent suffisamment pour compenser.
    
    \item \textbf{La value loss seule} : 0{,}0079 $\rightarrow$ 0{,}0096 (+21{,}3\%). Cette composante isolée augmente.
\end{itemize}

\subsection{Formule de la value loss}

\begin{equation}
L_{\text{value}} = \frac{1}{B}\sum_i (V(s_i) - r_i)^2
\end{equation}

C'est le MSE entre la prédiction de valeur $V(s)$ et la récompense réelle $r$.

\subsection{Pourquoi elle augmente}

Le \textbf{critique} (la tête de valeur du réseau) est censé prédire la récompense future. S'il diverge, c'est qu'il n'arrive plus à suivre la distribution des récompenses.

\textbf{Causes possibles :}
\begin{itemize}
    \item La politique change plus vite que le critique ne peut s'adapter
    \item Les phrases deviennent plus variées en récompense (plus difficile à prédire)
    \item Le learning rate est le même pour le critique et l'acteur, alors qu'il devrait être plus stable pour le critique
    \item La distribution des récompenses change (non-stationnarité)
\end{itemize}

\subsection{Pourquoi c'est anormal}

En principe, le critique devrait \textbf{converger} (value loss $\downarrow$), pas diverger. L'augmentation indique un problème de stabilité dans l'apprentissage.

\subsection{Masquage dans le plot global}

L'augmentation de 0{,}008 à 0{,}010 est minime en valeur absolue. Pondérée par 0{,}5 dans la loss totale, c'est +0{,}001. Cette contribution est noyée par la descente de la policy loss (-1{,}0) et de la contrastive loss (-0{,}9), donnant une descente globale visible dans le plot.

\section{Question 9 — Rôle de l'entropie en RL}

\subsection{Définition}

\textbf{L'entropie mesure le degré d'exploration de l'agent.}

\begin{itemize}
    \item Entropie haute : les actions sont variées et imprévisibles
    \item Entropie basse : l'agent fait toujours la même chose
\end{itemize}

\subsection{Formule dans votre implémentation}

L'entropie est un \textbf{proxy} calculé comme :
\begin{equation}
H = \overline{\sigma_{\text{dim}}(\hat{a})} + 0{,}5 \cdot \overline{\text{pdist}(\hat{a})}
\end{equation}

C'est la variance des prédictions + la diversité des prédictions dans le batch.

\subsection{Pourquoi on l'ajoute dans la loss}

Dans la loss finale : $L - 0{,}1 \cdot H$

On veut \textbf{maximiser} l'entropie (= encourager l'exploration). Donc on soustrait $H$ de la loss. Minimiser $L - 0{,}1 \cdot H$ revient à minimiser $L$ tout en maximisant $H$.

\subsection{Pourquoi c'est critique}

\textbf{Sans le bonus d'entropie, l'agent convergerait trop vite vers une seule stratégie} (toujours choisir le même type de prompt). Ce problème s'appelle le \textbf{policy collapse}.

L'entropie force l'agent à continuer d'explorer des alternatives, potentiellement meilleures.

\subsection{Exemples concrets}

\begin{itemize}
    \item \textbf{Sans bonus d'entropie} : L'agent apprend « prompt court = bon CER » et ne choisit plus que des prompts courts, ignorant que certains prompts longs pourraient être meilleurs.
    
    \item \textbf{Avec bonus d'entropie} : L'agent continue d'essayer des prompts de longueurs variées, même après avoir trouvé une stratégie qui marche bien.
\end{itemize}

\section{Question 10 — La descente active de la contrastive loss}

\subsection{Observation}

« Loss : Avg Loss encore plus élevée ($\sim$52{,}1, max 67{,}0) mais avec une tendance descendante nette (-6{,}79 entre premiers et derniers 10). La contrastive loss démarre à $\sim$198 et descend à $\sim$185, montrant un apprentissage actif des embeddings. »

\subsection{Contexte de pipeline\_261012 (filtered)}

Le modèle part d'un checkpoint précoce (itération 62). Les embeddings n'ont pas encore été bien calibrés — le réseau ne sait pas encore bien mapper les entrées vers l'espace des prompts.

Résultat : contrastive loss très élevée au départ (198).

\subsection{Ce que montre la descente 198 → 185}

\begin{enumerate}
    \item Le réseau apprend à produire des prédictions plus proches des bons prompts
    \item La projection input $\rightarrow$ espace 256D devient plus discriminante
    \item Le ratio $\frac{e^{\text{sim}(p, a^+)/\tau}}{\text{dénominateur}}$ augmente
    \item $-\log(\text{ratio})$ diminue $\rightarrow$ contrastive loss descend
\end{enumerate}

\subsection{Comparaison avec pipeline\_261011 (unfiltered)}

\begin{table}[h]
\centering
\begin{tabular}{lcc}
\toprule
\textbf{Métrique} & \textbf{Filtered (261012)} & \textbf{Unfiltered (261011)} \\
\midrule
Contrastive initiale & 198 & 193 \\
Contrastive finale & 185 & 190 \\
Descente & -13 (-6{,}6\%) & -3 (-1{,}6\%) \\
\bottomrule
\end{tabular}
\end{table}

Le filtered part de plus haut mais descend \textbf{plus vite par itération}. C'est un signe positif : le modèle est en phase d'apprentissage rapide.

\subsection{Pourquoi le filtrage aide}

Le filtrage élimine les phrases bruitées ou aberrantes. Cela donne un \textbf{signal de gradient plus propre} et plus exploitable :
\begin{itemize}
    \item Moins de variance dans les récompenses
    \item Pas de cas extrêmes qui tirent les paramètres dans des directions contradictoires
    \item La contrastive loss peut effectivement optimiser la séparation des embeddings
\end{itemize}

\section{Question 11 — Pourquoi la loss oscille dans les trois cas}

\subsection{Trois raisons fondamentales}

\begin{enumerate}
    \item \textbf{Batch size effectif = 1 pour la collecte d'expérience} : Chaque phrase génère une récompense unique. C'est comme estimer une moyenne en ne regardant qu'un échantillon à la fois — la variance est maximale. En supervised learning, un batch de 64 images lisse le gradient. Ici, même si le replay buffer permet de ré-entraîner sur 64 expériences, la récompense de chaque phrase varie énormément.
    
    \item \textbf{La récompense est stochastique et dépend de la phrase cible} : 
    \begin{itemize}
        \item « Bonjour » (1 mot) $\rightarrow$ CER probablement élevé $\rightarrow$ faible récompense
        \item « Le petit chat est mort sous la pluie d'automne » (9 mots) $\rightarrow$ CER probablement bon $\rightarrow$ forte récompense
    \end{itemize}
    Chaque phrase est un « environnement » différent.
    
    \item \textbf{$\varepsilon$-greedy exploration} : Avec $\varepsilon > 0$, l'agent choisit parfois un prompt aléatoire au lieu du meilleur. Ce prompt aléatoire peut être catastrophique, créant un pic de loss.
\end{enumerate}

\subsection{Pourquoi pipeline\_261012 est le plus oscillant}

\begin{table}[h]
\centering
\begin{tabular}{lccc}
\toprule
\textbf{Métrique} & \textbf{Filtered (261012)} & \textbf{Unfiltered (261011)} & \textbf{Baseline (260326)} \\
\midrule
Early volatility & 1{,}07 & 0{,}40 & 0{,}14 \\
Late volatility & 0{,}89 & 0{,}33 & 0{,}07 \\
\bottomrule
\end{tabular}
\caption{Volatilité de la loss (écart-type glissant)}
\end{table}

\textbf{Raisons :}
\begin{itemize}
    \item Le modèle part d'un état précoce (value loss initiale 0{,}048, soit \textbf{10× plus} que l'unfiltered)
    \item Les gradients sont donc beaucoup plus grands car l'erreur est plus grande
    \item La contrastive loss commence à 198--225 et subit de fortes variations (std = 5{,}37 vs 1{,}58 pour l'unfiltered)
    \item Chaque batch du replay buffer a une composition très différente
\end{itemize}

\textbf{Analogie :} C'est comme une balle lancée du haut d'une colline (gros mouvements) vs une balle près du fond de la vallée (petits ajustements). Plus on est loin de la convergence, plus les oscillations sont fortes.

\section{Question 12 — Apprentissage plus rapide dans le filtered}

\subsection{Observation}

« Observation clé : la version filtrée montre un apprentissage beaucoup plus rapide par itération. La value loss passe de 0{,}048 à 0{,}007 (réduction de 86\%), ce qui est excellent. La policy loss se creuse de -2{,}2 à -5{,}0, montrant que l'agent apprend efficacement à maximiser la récompense. »

\subsection{Value loss : 0{,}048 → 0{,}007 (réduction 86\%)}

\textbf{La value loss initiale haute (0{,}048) signifie que le critique faisait de grosses erreurs de prédiction au départ.}

\begin{equation}
L_{\text{value}} = \frac{1}{B}\sum_i (V(s_i) - r_i)^2
\end{equation}

Le gradient de cette loss est proportionnel à l'erreur : $(V(s) - r)$.

\begin{itemize}
    \item Quand l'erreur est grande, le gradient est grand $\rightarrow$ l'apprentissage est rapide
    \item C'est la propriété naturelle de la MSE loss
\end{itemize}

Le passage à 0{,}007 montre que le critique a appris à prédire précisément les récompenses.

\textbf{Comparaison :} L'unfiltered commence déjà à 0{,}008 (le critique est déjà assez bon), donc il a peu de marge d'amélioration, et en fait il diverge légèrement (+21{,}3\%).

\subsection{Policy loss : -2{,}2 → -5{,}0}

La policy loss qui se creuse signifie que l'agent produit des prédictions de plus en plus alignées avec les bonnes actions.

\begin{itemize}
    \item -2{,}2 au départ : la cosine similarity pondérée est médiocre
    \item -5{,}0 à la fin : l'agent prédit des vecteurs très proches des bons prompts, pondérés par des avantages bien estimés (grâce au critique amélioré)
\end{itemize}

\subsection{Pourquoi le filtrage aide}

Le filtrage élimine les échantillons bruités (phrases aberrantes, CER $>$ 1, récompenses 0{,}0).

\textbf{Cela donne un signal de gradient plus propre :}
\begin{enumerate}
    \item Moins de variance dans les récompenses $\rightarrow$ avantages mieux estimés
    \item Pas de cas extrêmes qui tirent les paramètres dans des directions contradictoires
    \item Le critique a une distribution de récompenses plus stable à prédire
    \item Les mises à jour sont cohérentes (pas de bruit contradictoire)
\end{enumerate}

\subsection{Tableau comparatif}

\begin{table}[h]
\centering
\begin{tabular}{lcc}
\toprule
\textbf{Métrique} & \textbf{Filtered (261012)} & \textbf{Unfiltered (261011)} \\
\midrule
Value loss initiale & 0{,}048 & 0{,}008 \\
Value loss finale & 0{,}007 & 0{,}010 \\
\textbf{Réduction} & \textbf{-86{,}2\%} & \textbf{+21{,}3\%} (diverge) \\
\midrule
Policy loss initiale & -2{,}2 & -4{,}2 \\
Policy loss finale & -5{,}0 & -5{,}7 \\
\textbf{Amélioration} & \textbf{-2{,}8} & \textbf{-1{,}5} \\
\bottomrule
\end{tabular}
\caption{Dynamique d'apprentissage : filtered vs unfiltered}
\end{table}

\section{Question 13 — « Le signal InfoNCE est noyé dans le bruit »}

\subsection{Contexte}

« Les embeddings ne séparent pas les candidats. Le signal InfoNCE est noyé dans le bruit. »

\subsection{Explication mathématique}

Rappel de la formule InfoNCE :
\begin{equation}
p(\text{positif}) = \frac{e^{\text{sim}(p, a^+)/\tau}}{e^{\text{sim}(p, a^+)/\tau} + \sum_j e^{\text{sim}(p, a^-_j)/\tau}}
\end{equation}

\textbf{Exemple numérique :} positif sim = 0{,}80, négatifs sim = 0{,}78, 0{,}79, $\ldots$ (presque identiques)

Avec $\tau = 0{,}1$ et 63 négatifs :
\begin{align}
p &\approx \frac{e^{0{,}8/0{,}1}}{e^{0{,}8/0{,}1} + 63 \times e^{0{,}79/0{,}1}} \\
&= \frac{e^{8}}{e^{8} + 63 \times e^{7{,}9}} \\
&\approx \frac{2981}{2981 + 63 \times 2697} \\
&\approx \frac{2981}{172{,}892} \\
&\approx 0{,}017
\end{align}

La loss : $-\log(0{,}017) \approx 4{,}07$

\textbf{Cas aléatoire pur :} $-\log(1/64) = 4{,}16$

\subsection{Interprétation}

La loss InfoNCE est \textbf{quasiment identique} au cas où l'agent choisirait aléatoirement parmi les 64 candidats.

\textbf{« Noyé dans le bruit » signifie :}
\begin{itemize}
    \item Le signal informatif (différence entre positif et négatifs) est si faible qu'il disparaît
    \item La loss ne contient presque aucune information utile pour l'apprentissage
    \item Le gradient résultant pointe dans une direction presque aléatoire
    \item L'agent n'apprend rien d'utile via cette composante
\end{itemize}

\subsection{Cas du pipeline v2}

Avec une contrastive loss de 192, on est bien au-delà du cas aléatoire :
\begin{equation}
192 \approx -\log(10^{-83})
\end{equation}

Le modèle assigne une probabilité quasi-nulle au bon prompt. Il ne distingue \textbf{littéralement pas} le bon prompt des mauvais.

\section{Question 14 — Policy collapse}

\subsection{Définition}

\textbf{Policy collapse} = l'agent a trouvé \textbf{une seule stratégie qui marche assez bien} et s'y accroche, cessant d'explorer des alternatives potentiellement meilleures.

\subsection{Concrètement dans votre cas}

L'agent pourrait apprendre que « les prompts courts donnent un bon CER ». 

\begin{itemize}
    \item Il se fixe sur cette stratégie : toujours sélectionner le prompt le plus court parmi les candidats
    \item Son entropie baisse (0{,}720 $\rightarrow$ 0{,}689)
    \item Il ne diversifie plus ses choix
\end{itemize}

\subsection{Le problème}

Peut-être qu'un prompt de 12 mots avec un contenu sémantiquement proche de la cible donnerait un CER encore meilleur, mais l'agent ne l'essaie jamais car il a convergé vers « court = bon ».

\textbf{La stratégie actuelle est sous-optimale, mais l'agent ne le saura jamais.}

\subsection{Indices de policy collapse}

\begin{enumerate}
    \item Entropie qui descend continuellement
    \item Toujours les mêmes types de prompts sélectionnés
    \item La reward stagne à un plateau (0{,}64) au lieu de continuer à monter
    \item L'agent ignore certaines régions de l'espace des actions
    \item Distribution des actions très piquée (faible diversité)
\end{enumerate}

\subsection{Mécanisme mathématique}

\begin{equation}
\pi(a|s) = \text{softmax}(\text{logits})
\end{equation}

Sans régularisation, la politique devient de plus en plus déterministe :
\begin{itemize}
    \item $\pi(a_{\text{préféré}}|s) \rightarrow 1$
    \item $\pi(a_{\text{autres}}|s) \rightarrow 0$
\end{itemize}

Plus de gradient exploitable pour améliorer les alternatives.

\subsection{Comment l'éviter}

\begin{enumerate}
    \item \textbf{Bonus d'entropie} : $-0{,}1 \cdot H$ dans la loss (déjà implémenté)
    \item \textbf{Maintenir $\varepsilon > 0$} : forcer l'exploration aléatoire
    \item \textbf{Augmenter le coefficient d'entropie} si le collapse est observé (ex: passer de 0{,}1 à 0{,}2)
    \item \textbf{Replay buffer diversifié} : s'assurer que le buffer contient des expériences variées
    \item \textbf{Curriculum learning} : augmenter progressivement la difficulté des phrases
\end{enumerate}

\subsection{Observation dans vos données}

Le baseline montre des signes de début de collapse :
\begin{itemize}
    \item Entropie : 0{,}748 $\rightarrow$ 0{,}740 (-1{,}1\%)
    \item Reward stagne puis régresse légèrement
    \item Prompt selection : 59{,}5\% prompts courts vs 16{,}7\% prompts longs (ratio 3{,}57×)
\end{itemize}

Le filtrage pourrait aggraver le problème si on filtre toujours les mêmes types de phrases difficiles, réduisant encore la diversité.

\newpage
\part{Questions supplémentaires}

\section{Question 15 — Explication terme par terme de la policy loss}

\subsection{Formule complète}

\begin{equation}
L_{\text{policy}} = -\frac{1}{B}\sum_{i=1}^{B} \log\!\big(\cos(\hat{a}_i, a_i)\big) \cdot \hat{A}_i
\end{equation}

\subsection{Définition de chaque variable}

\begin{description}[style=nextline, leftmargin=2cm]
    \item[$B$] La taille du batch. C'est le nombre d'expériences (phrase + prompt + récompense) sur lesquelles on calcule la loss en une seule passe. La somme est moyennée par $1/B$ pour que la loss ne dépende pas de la taille du batch.
    
    \item[$\hat{a}_i$] La \textbf{prédiction de l'agent} : un vecteur 256D produit par le réseau de neurones. C'est l'embedding que l'agent propose comme « action idéale ». Il représente la direction dans l'espace latent vers laquelle l'agent pense que le meilleur prompt se trouve.
    
    \item[$a_i$] L'\textbf{action réellement exécutée} : l'embedding 256D du prompt qui a été effectivement sélectionné (via FAISS, en cherchant le voisin le plus proche de $\hat{a}_i$). C'est le vecteur du prompt candidat qui a été utilisé pour générer l'audio et obtenir le CER.
    
    \item[$\cos(\hat{a}_i, a_i)$] La \textbf{similarité cosinus} entre la prédiction et l'action exécutée :
    \begin{equation}
    \cos(\hat{a}, a) = \frac{\hat{a} \cdot a}{\|\hat{a}\| \cdot \|a\|}
    \end{equation}
    Elle mesure à quel point la prédiction de l'agent est alignée avec le prompt qu'il a choisi. Valeur dans $[0, 1]$ (après normalisation L2). Plus c'est proche de 1, plus la prédiction est précise.
    
    \item[$\log\!\big(\cos(\hat{a}_i, a_i)\big)$] Le \textbf{logarithme} de la similarité cosinus. Puisque $\cos \in (0, 1]$, le log est toujours $\leq 0$ :
    \begin{itemize}
        \item $\cos = 1 \Rightarrow \log(1) = 0$ (prédiction parfaite)
        \item $\cos = 0{,}5 \Rightarrow \log(0{,}5) \approx -0{,}69$ (prédiction médiocre)
        \item $\cos \to 0 \Rightarrow \log \to -\infty$ (prédiction catastrophique)
    \end{itemize}
    Le log amplifie les pénalités pour les mauvaises prédictions (effet similaire à la cross-entropy en classification).
    
    \item[$\hat{A}_i$] L'\textbf{avantage normalisé}. Il indique si l'action prise était meilleure ou pire que la moyenne.
    
    \item[$-$] Le \textbf{signe négatif} devant toute l'expression. Il transforme la minimisation de la loss en maximisation de la récompense attendue. Sans ce signe, l'optimiseur minimiserait la récompense au lieu de la maximiser.
\end{description}

\subsection{Le rôle de la similarité cosinus}

La similarité cosinus joue le rôle de la \textbf{probabilité d'action} dans l'algorithme REINFORCE classique.

En REINFORCE standard :
\begin{equation}
L = -\log \pi(a|s) \cdot \hat{A}
\end{equation}

Dans notre cas, l'action n'est pas un choix discret (choisir parmi $N$ options), mais un \textbf{vecteur continu} dans $\mathbb{R}^{256}$. On remplace donc $\pi(a|s)$ par $\cos(\hat{a}, a)$ : plus la prédiction est proche de l'action choisie, plus la « probabilité » est élevée.

\subsection{Le sens du signe négatif}

Le signe négatif est essentiel car on utilise un \textbf{optimiseur qui minimise} (descente de gradient).

\textbf{Objectif réel :} maximiser $\mathbb{E}[\log(\cos) \cdot \hat{A}]$

\textbf{Ce qu'on minimise :} $-\mathbb{E}[\log(\cos) \cdot \hat{A}]$

Minimiser le négatif d'une quantité = maximiser cette quantité.

\textbf{Exemple :}
\begin{itemize}
    \item $\cos = 0{,}9$, $\hat{A} = +2$ (bonne action) : $L = -\log(0{,}9) \times 2 = -(-0{,}105) \times 2 = -0{,}21$. On veut que la loss descende $\rightarrow$ on renforce cette action.
    \item $\cos = 0{,}3$, $\hat{A} = -1$ (mauvaise action) : $L = -\log(0{,}3) \times (-1) = -(-1{,}2) \times (-1) = -1{,}2$. On veut que la loss monte $\rightarrow$ on évite cette action.
\end{itemize}

\subsection{L'avantage $\hat{A}_i$ : définition et calcul}

\textbf{Définition :} l'avantage mesure si l'action prise a produit une récompense \textbf{meilleure ou pire que la moyenne}.

\textbf{Calcul dans le code :}
\begin{equation}
\hat{A}_i = \frac{r_i - \bar{r}}{\sigma_r + 10^{-8}}
\end{equation}

où :
\begin{itemize}
    \item $r_i$ : la récompense obtenue pour cette action (= $1 - \text{CER}$)
    \item $\bar{r}$ : la récompense moyenne dans le replay buffer (baseline)
    \item $\sigma_r$ : l'écart-type des récompenses dans le buffer
    \item $10^{-8}$ : constante pour éviter la division par zéro
\end{itemize}

\textbf{Avantage positif ($\hat{A} > 0$) :} la récompense de cette action est \textbf{supérieure à la moyenne}. L'agent doit \textbf{renforcer} cette action (produire des prédictions plus proches de ce prompt à l'avenir).

\textbf{Avantage négatif ($\hat{A} < 0$) :} la récompense est \textbf{en dessous de la moyenne}. L'agent doit \textbf{s'éloigner} de cette action (produire des prédictions différentes de ce prompt).

\textbf{Avantage nul ($\hat{A} = 0$) :} la récompense est exactement la moyenne. Pas de mise à jour.

\subsection{Exemple complet}

\begin{table}[h]
\centering
\begin{tabular}{cccccc}
\toprule
$r_i$ & $\bar{r}$ & $\hat{A}_i$ & $\cos$ & $\log(\cos)$ & $L_i$ \\
\midrule
0{,}9 & 0{,}6 & +1{,}5 & 0{,}95 & -0{,}05 & $-(-0{,}05)(1{,}5) = -0{,}075$ \\
0{,}3 & 0{,}6 & -1{,}5 & 0{,}80 & -0{,}22 & $-(-0{,}22)(-1{,}5) = +0{,}33$ \\
0{,}6 & 0{,}6 & 0{,}0 & 0{,}70 & -0{,}36 & $-(-0{,}36)(0) = 0$ \\
\bottomrule
\end{tabular}
\caption{Effet de l'avantage sur la policy loss}
\end{table}

\begin{itemize}
    \item Ligne 1 : bonne action ($\hat{A} > 0$), bonne prédiction ($\cos$ élevé) $\rightarrow$ loss négative, renforce l'action.
    \item Ligne 2 : mauvaise action ($\hat{A} < 0$) $\rightarrow$ loss positive, pénalise l'action.
    \item Ligne 3 : action moyenne ($\hat{A} = 0$) $\rightarrow$ pas de contribution.
\end{itemize}

\section{Question 16 — Explication terme par terme de la contrastive loss (InfoNCE)}

\subsection{Formule complète}

\begin{equation}
L_{\text{InfoNCE}} = -\log \frac{e^{\text{sim}(p, a^+)/\tau}}{e^{\text{sim}(p, a^+)/\tau} + \sum_{j=1}^{N} e^{\text{sim}(p, a_j^-)/\tau}}
\end{equation}

\subsection{Définition de chaque terme}

\begin{description}[style=nextline, leftmargin=2cm]
    \item[$p$] La \textbf{prédiction} de l'agent : le vecteur 256D produit par le réseau. C'est le point dans l'espace latent où l'agent « pointe » pour indiquer le type de prompt qu'il veut.
    
    \item[$a^+$] Le \textbf{prompt positif} : l'embedding du prompt qui a donné la meilleure récompense (meilleur CER). C'est l'action « correcte » — celle que l'agent devrait apprendre à prédire. Dans le code, c'est l'action dont la récompense est supérieure à la baseline.
    
    \item[$a_j^-$] Les \textbf{prompts négatifs} : les embeddings des prompts qui ont donné de mauvaises récompenses (CER élevé). Il y en a $N$ dans le batch. Ce sont les « distracteurs » dont l'agent doit apprendre à se distinguer.
    
    \item[$\text{sim}(p, a)$] La \textbf{similarité} entre deux vecteurs. Dans le code, c'est le produit scalaire après normalisation L2 (équivalent à la cosine similarity) :
    \begin{equation}
    \text{sim}(p, a) = \frac{p}{\|p\|} \cdot \frac{a}{\|a\|} = \cos(p, a)
    \end{equation}
    Valeur dans $[-1, 1]$. Plus c'est proche de 1, plus les vecteurs pointent dans la même direction.
    
    \item[$\tau$] La \textbf{température}. C'est un hyperparamètre qui contrôle la « netteté » de la distribution softmax :
    \begin{itemize}
        \item $\tau$ petit (0{,}05) : distribution très piquée, le modèle est très discriminant. Petites différences de similarité $\rightarrow$ grandes différences de probabilité.
        \item $\tau$ grand (0{,}5) : distribution plate, le modèle traite les candidats de manière plus uniforme.
    \end{itemize}
    Dans le code : $\tau = 0{,}1$ (fixe) ou adaptatif entre 0{,}05 et 0{,}5.
    
    \item[$e^{\text{sim}/\tau}$] L'\textbf{exponentielle} de la similarité divisée par $\tau$. C'est la transformation softmax. Elle rend les scores de similarité comparables en les transformant en « pseudo-probabilités ».
    
    \item[$\frac{e^{\text{sim}(p, a^+)/\tau}}{\text{dénominateur}}$] Le \textbf{ratio softmax}. Il représente la probabilité que le modèle assigne au prompt positif parmi tous les candidats (positif + négatifs). Plus ce ratio est proche de 1, plus le modèle est confiant que $a^+$ est le bon prompt.
    
    \item[$-\log(\cdot)$] Le \textbf{logarithme négatif}. Il transforme la probabilité en loss :
    \begin{itemize}
        \item Ratio $= 1$ (confiance parfaite) $\Rightarrow$ loss $= 0$
        \item Ratio $= 0{,}5$ $\Rightarrow$ loss $= 0{,}69$
        \item Ratio $\to 0$ (incertitude totale) $\Rightarrow$ loss $\to +\infty$
    \end{itemize}
\end{description}

\subsection{Rôle des exemples positifs et négatifs}

\textbf{Le positif ($a^+$)} sert de « cible » : l'agent doit produire des prédictions proches de ce prompt.

\textbf{Les négatifs ($a_j^-$)} servent de « contraste » : l'agent doit produire des prédictions éloignées de ces prompts. Plus il y a de négatifs, plus la tâche est difficile (le dénominateur est plus grand), mais plus le signal d'apprentissage est riche.

\textbf{Analogie :} c'est comme un QCM à 64 choix. $a^+$ est la bonne réponse, les 63 $a_j^-$ sont les mauvaises réponses. L'InfoNCE pousse le modèle à donner la probabilité la plus haute à la bonne réponse.

\section{Question 17 — Embeddings bien ou mal séparés dans l'espace latent}

\subsection{L'espace latent}

L'espace latent est un espace vectoriel de dimension 256 ($\mathbb{R}^{256}$) où chaque prompt est représenté par un point (un vecteur). Deux prompts qui produisent des résultats similaires devraient avoir des vecteurs proches, et deux prompts très différents devraient avoir des vecteurs éloignés.

\subsection{Bien séparés — Exemple concret}

Imaginons un espace simplifié en 2D pour visualiser.

\textbf{Bien séparés :}
\begin{itemize}
    \item Prompt A (« bonjour monsieur ») $\rightarrow$ vecteur $(0{,}9,\ 0{,}1)$
    \item Prompt B (« la pluie tombe ») $\rightarrow$ vecteur $(-0{,}3,\ 0{,}8)$
    \item Prompt C (« il fait beau ») $\rightarrow$ vecteur $(-0{,}7,\ -0{,}5)$
\end{itemize}

Chaque prompt occupe une zone distincte de l'espace. Si l'agent prédit un vecteur proche de $(0{,}9,\ 0{,}1)$, on sait sans ambiguïté qu'il veut le prompt A. La cosine similarity avec A sera ~0{,}95, et avec B ou C elle sera ~0{,}2.

Le ratio InfoNCE :
\begin{equation}
\frac{e^{0{,}95/0{,}1}}{e^{0{,}95/0{,}1} + e^{0{,}2/0{,}1} + e^{0{,}2/0{,}1}} = \frac{e^{9{,}5}}{e^{9{,}5} + 2 \times e^{2}} \approx \frac{13360}{13360 + 14{,}8} \approx 0{,}999
\end{equation}

Loss $= -\log(0{,}999) \approx 0{,}001$ — \textbf{très faible}.

\subsection{Mal séparés — Exemple concret}

\textbf{Mal séparés :}
\begin{itemize}
    \item Prompt A $\rightarrow$ vecteur $(0{,}71,\ 0{,}70)$
    \item Prompt B $\rightarrow$ vecteur $(0{,}72,\ 0{,}69)$
    \item Prompt C $\rightarrow$ vecteur $(0{,}70,\ 0{,}71)$
\end{itemize}

Tous les prompts sont regroupés dans la même zone. La prédiction de l'agent aura une cosine similarity presque identique avec les trois :
\begin{equation}
\cos(\hat{a}, A) \approx 0{,}998, \quad \cos(\hat{a}, B) \approx 0{,}997, \quad \cos(\hat{a}, C) \approx 0{,}996
\end{equation}

Le ratio InfoNCE :
\begin{equation}
\frac{e^{0{,}998/0{,}1}}{e^{0{,}998/0{,}1} + e^{0{,}997/0{,}1} + e^{0{,}996/0{,}1}} \approx \frac{e^{9{,}98}}{e^{9{,}98} + e^{9{,}97} + e^{9{,}96}} \approx \frac{1}{3} = 0{,}33
\end{equation}

Loss $= -\log(0{,}33) \approx 1{,}1$ — \textbf{élevée} car le modèle ne peut pas distinguer le bon prompt des mauvais.

\subsection{Résumé visuel}

\begin{table}[h]
\centering
\begin{tabular}{lcc}
\toprule
\textbf{Situation} & \textbf{InfoNCE ratio} & \textbf{Loss} \\
\midrule
Bien séparés & $\approx 0{,}999$ & $\approx 0{,}001$ \\
Mal séparés & $\approx 0{,}33$ & $\approx 1{,}1$ \\
Cas aléatoire (64 candidats) & $\approx 0{,}016$ & $\approx 4{,}16$ \\
Pipeline v2 (192) & $\approx 10^{-83}$ & $\approx 192$ \\
\bottomrule
\end{tabular}
\caption{Lien entre séparation des embeddings et contrastive loss}
\end{table}

\section{Question 18 — La value loss : définition et fonctionnement}

\subsection{Définition}

La \textbf{value loss} entraîne le \textbf{critique} — une partie du réseau qui apprend à prédire quelle récompense l'agent obtiendra dans un état donné. Sa formule :

\begin{equation}
L_{\text{value}} = \frac{1}{B}\sum_{i=1}^{B} \big(V(s_i) - r_i\big)^2
\end{equation}

où :
\begin{itemize}
    \item $V(s_i)$ : la valeur prédite par le critique pour l'état $s_i$ (l'embedding de la phrase cible)
    \item $r_i$ : la récompense réellement obtenue après avoir agi
\end{itemize}

\subsection{« On ne connaît pas la récompense future » — Explication}

C'est vrai qu'au moment de la \textbf{prédiction}, le critique ne connaît pas la récompense réelle. Il fait une \textbf{estimation a priori}.

Mais la value loss n'est pas calculée \textit{au moment de la prédiction}. Elle est calculée \textbf{après coup}, lors de la phase d'entraînement :

\begin{enumerate}
    \item \textbf{Phase d'action :} L'agent reçoit une phrase cible $s_i$. Le critique prédit $V(s_i) = 0{,}7$ (« je pense que la récompense sera 0{,}7 »). L'agent sélectionne un prompt, génère l'audio, et obtient le CER.
    
    \item \textbf{Phase de retour :} On calcule la récompense réelle $r_i = 1 - \text{CER} = 0{,}6$.
    
    \item \textbf{Phase d'entraînement :} On compare la prédiction à la réalité :
    \begin{equation}
    L_{\text{value}} = (V(s_i) - r_i)^2 = (0{,}7 - 0{,}6)^2 = 0{,}01
    \end{equation}
    Le gradient de cette erreur est rétropropagé pour que le critique prédise mieux la prochaine fois.
\end{enumerate}

\subsection{Ce que représente la « valeur cible »}

La \textbf{valeur cible} est simplement la récompense $r_i$ que l'on a \textbf{effectivement observée}. Ce n'est pas une prédiction — c'est un fait mesuré après l'exécution de l'action.

\textbf{Analogie :} un météorologue prédit « 20°C demain ». Le lendemain, on mesure 18°C. L'erreur $(20 - 18)^2 = 4$ sert à calibrer le modèle de prévision. La mesure réelle (18°C) est la « valeur cible ».

C'est exactement le même principe : le critique fait une prédiction \textbf{avant}, et on corrige \textbf{après} avec la vraie valeur.

\section{Question 19 — Comment le modèle prédit-il les récompenses futures ?}

\subsection{Architecture du critique}

Le critique est une \textbf{tête du réseau de neurones} (un MLP séparé) qui partage les couches initiales avec l'acteur :

\begin{enumerate}
    \item L'embedding de la phrase cible (1024D, SONAR) entre dans le réseau
    \item Les couches partagées (projection + blocs résiduels) transforment cet input
    \item \textbf{L'acteur} : produit un vecteur 256D (la prédiction d'action)
    \item \textbf{Le critique} : produit un scalaire (la valeur prédite $V(s)$)
\end{enumerate}

\subsection{Comment il apprend}

Le critique apprend par \textbf{régression supervisée sur les récompenses passées}. À chaque itération :

\begin{enumerate}
    \item Il reçoit l'état $s$ (embedding de la phrase cible)
    \item Il prédit $V(s)$ (« cette phrase devrait donner une récompense de X »)
    \item On observe la récompense réelle $r$
    \item On calcule l'erreur : $(V(s) - r)^2$
    \item On met à jour les poids pour réduire cette erreur
\end{enumerate}

Après des centaines d'itérations, le critique a vu suffisamment de paires (phrase, récompense) pour \textbf{apprendre des patterns statistiques}. Par exemple :
\begin{itemize}
    \item Phrases longues $\rightarrow$ récompenses généralement élevées ($\sim$0{,}73)
    \item Phrases très courtes $\rightarrow$ récompenses généralement faibles ($\sim$0{,}57)
    \item Phrases avec vocabulaire courant $\rightarrow$ meilleures récompenses
\end{itemize}

\subsection{Pourquoi il peut diverger}

« Le critique diverge » signifie que ses prédictions deviennent \textbf{de moins en moins précises} au lieu de s'améliorer.

\textbf{Causes :}
\begin{enumerate}
    \item \textbf{Non-stationnarité :} La politique de l'acteur change au cours du temps. Cela modifie la distribution des récompenses. Le critique apprend à prédire les récompenses d'une politique qui n'existe plus.
    
    \item \textbf{Learning rate inadapté :} Si le learning rate est trop élevé pour le critique, il « oublie » ses anciennes connaissances à chaque mise à jour. Idéalement, le critique devrait avoir un learning rate plus faible que l'acteur.
    
    \item \textbf{Replay buffer périmé :} Les expériences anciennes dans le buffer correspondent à une politique différente. Le critique apprend sur des données « périmées ».
\end{enumerate}

\subsection{Conséquence de la divergence}

Si le critique est imprécis, l'avantage $\hat{A} = r - V(s)$ est mal estimé. Les mises à jour de la politique sont alors bruitées : l'agent renforce parfois de mauvaises actions et pénalise de bonnes actions. C'est un cercle vicieux qui déstabilise tout l'apprentissage.

\section{Question 20 — Acteur et critique : rôles respectifs}

\subsection{Architecture Actor-Critic}

Le pipeline utilise une architecture \textbf{Actor-Critic}, où un seul réseau de neurones possède deux « têtes » distinctes :

\subsection{L'acteur (Actor)}

\begin{itemize}
    \item \textbf{Qui :} la tête du réseau qui produit l'\textbf{action} (vecteur 256D)
    \item \textbf{Entrée :} l'embedding SONAR de la phrase cible (1024D)
    \item \textbf{Sortie :} un vecteur 256D normalisé qui pointe vers le meilleur prompt estimé
    \item \textbf{Rôle :} décider quelle action prendre — quel prompt sélectionner
    \item \textbf{Entraîné par :} la policy loss ($L_{\text{policy}}$) et la contrastive loss ($L_{\text{InfoNCE}}$)
    \item \textbf{Objectif :} maximiser la récompense attendue
\end{itemize}

\subsection{Le critique (Critic)}

\begin{itemize}
    \item \textbf{Qui :} la tête du réseau qui produit une \textbf{estimation de valeur} (un scalaire)
    \item \textbf{Entrée :} les mêmes features internes que l'acteur (couches partagées)
    \item \textbf{Sortie :} un nombre réel $V(s) \in [0, 1]$ qui estime la récompense que l'agent obtiendra
    \item \textbf{Rôle :} évaluer la qualité de l'état actuel (la phrase cible) — « est-ce qu'on va bien s'en sortir avec cette phrase ? »
    \item \textbf{Entraîné par :} la value loss ($L_{\text{value}}$)
    \item \textbf{Objectif :} prédire précisément la récompense future
\end{itemize}

\subsection{Interaction entre les deux}

Le critique \textbf{aide l'acteur} en fournissant la baseline pour calculer l'avantage :

\begin{equation}
\hat{A} = r - V(s) \quad \text{(simplifié, avant normalisation)}
\end{equation}

\begin{itemize}
    \item Si $r > V(s)$ : l'action était \textbf{meilleure que prévu} $\rightarrow$ renforcer
    \item Si $r < V(s)$ : l'action était \textbf{pire que prévu} $\rightarrow$ pénaliser
    \item Si $r = V(s)$ : l'action était \textbf{conforme aux attentes} $\rightarrow$ pas de changement
\end{itemize}

\textbf{Analogie :} l'acteur est un joueur de football, le critique est l'entraîneur. Le joueur (acteur) décide de l'action (tirer, passer). L'entraîneur (critique) évalue « est-ce que cette action a été meilleure ou pire que ce que j'attendais ? ». Le joueur s'améliore en écoutant l'évaluation de l'entraîneur.

\subsection{Pourquoi cette architecture ?}

\textbf{Sans le critique} (REINFORCE pur), la baseline serait simplement la moyenne des récompenses. Cela fonctionne, mais l'estimation de l'avantage est bruitée.

\textbf{Avec le critique} (Actor-Critic), la baseline est \textbf{conditionnelle à l'état} : pour une phrase facile, $V(s)$ sera élevé ; pour une phrase difficile, $V(s)$ sera bas. L'avantage est donc mieux calibré, ce qui réduit la variance des gradients et accélère l'apprentissage.

\section{Question 21 — Augmenter la dimension de l'espace latent (100D → 256D)}

\subsection{En théorie, oui — mais pas toujours en pratique}

Augmenter la dimension de l'espace latent donne \textbf{plus de capacité de représentation}. Mais comme souvent en machine learning, il y a un compromis.

\subsection{Avantages de 256D vs 100D}

\begin{enumerate}
    \item \textbf{Plus de capacité de séparation :} Dans un espace 256D, il y a plus de directions orthogonales disponibles pour séparer les prompts. En 100D, on a 100 axes indépendants. En 256D, on en a 256. Cela permet de mieux distinguer les prompts proches.
    
    \item \textbf{Meilleure expressivité :} Le réseau peut encoder plus de nuances dans les relations entre prompts. Des prompts qui étaient confondus en 100D (car les dimensions étaient saturées) peuvent être séparés en 256D.
    
    \item \textbf{Conservation de l'information :} L'embedding SONAR source est en 1024D. Projeter de 1024D vers 100D compresse beaucoup plus (perte de 90{,}2\% des dimensions) que vers 256D (perte de 75\%). En 256D, on conserve davantage l'information originale.
    
    \item \textbf{Meilleure contrastive loss potentielle :} Avec plus de dimensions, les vecteurs ont plus de « place » pour s'éloigner les uns des autres, ce qui facilite la tâche de l'InfoNCE.
\end{enumerate}

\subsection{Risques de 256D}

\begin{enumerate}
    \item \textbf{Curse of dimensionality :} En haute dimension, tous les vecteurs tendent à devenir \textbf{équidistants}. La cosine similarity entre deux vecteurs aléatoires en 256D se concentre autour de 0 (distribution $\mathcal{N}(0, 1/256)$). Cela peut rendre la discrimination plus difficile, pas plus facile.
    
    \item \textbf{Overfitting :} Plus de paramètres dans les couches de projection (1024 $\times$ 256 = 262{,}144 paramètres vs 1024 $\times$ 100 = 102{,}400). Si les données d'entraînement sont limitées, le réseau peut mémoriser au lieu de généraliser.
    
    \item \textbf{FAISS moins efficace :} La recherche de plus proches voisins dans un espace 256D est plus coûteuse et potentiellement moins précise qu'en 100D (le nombre d'échantillons nécessaires pour couvrir l'espace croît exponentiellement avec la dimension).
    
    \item \textbf{Dimensions inutiles :} Si la structure intrinsèque des données vit dans un sous-espace de 100D, les 156 dimensions supplémentaires ne contiennent que du bruit. Le réseau doit apprendre à les ignorer, ce qui gaspille de la capacité.
\end{enumerate}

\subsection{En pratique}

La question clé est : \textbf{la complexité intrinsèque des données justifie-t-elle 256D ?}

\begin{itemize}
    \item Si on a 63{,}000 prompts très diversifiés $\rightarrow$ 256D peut aider
    \item Si les prompts sont assez similaires (même domaine) $\rightarrow$ 100D peut suffire
    \item Il faudrait comparer les performances (CER, reward) sur les deux dimensions pour conclure
\end{itemize}

\textbf{Conclusion :} En théorie, 256D offre plus de potentiel, mais le gain n'est pas garanti. Il dépend de la diversité des données et de la quantité d'entraînement disponible.

\section{Question 22 — Policy loss de -5 = « prédictions très proches des bons prompts »}

\subsection{L'apparente contradiction}

La valeur -5 semble « mauvaise » car elle est loin de 0. Mais en réalité, dans le contexte de la policy loss RL, \textbf{plus c'est négatif, mieux c'est}.

\subsection{Décomposition}

Rappel de la formule :
\begin{equation}
L_{\text{policy}} = -\log(\cos(\hat{a}, a)) \cdot \hat{A}
\end{equation}

Pour obtenir $L_{\text{policy}} = -5$, il faut que le produit $\log(\cos) \cdot \hat{A}$ soit positif et grand. Cela se produit quand :

\begin{enumerate}
    \item $\cos(\hat{a}, a)$ est \textbf{proche de 1} (bonne prédiction). Par exemple $\cos = 0{,}95$.
    \item $\hat{A}$ est \textbf{nettement positif} (récompense très supérieure à la moyenne). Par exemple $\hat{A} = +2$.
\end{enumerate}

\textbf{Calcul :}
\begin{equation}
L = -\log(0{,}95) \times 2 = -(-0{,}05) \times 2 = -0{,}10 \quad \text{(pour une paire)}
\end{equation}

Pour obtenir un mean de $-5$, cela signifie que \textbf{beaucoup de paires} dans le batch contribuent avec des valeurs négatives, c'est-à-dire que la prédiction est systématiquement alignée avec des actions à haute récompense.

\subsection{Évolution au cours du training}

\begin{table}[h]
\centering
\begin{tabular}{ccl}
\toprule
\textbf{Policy loss} & \textbf{Signification} & \textbf{Qualité} \\
\midrule
$-2{,}2$ & Prédictions faiblement alignées & Début d'apprentissage \\
$-4{,}0$ & Prédictions moyennement alignées & En progression \\
$-5{,}0$ & Prédictions bien alignées & Bonne convergence \\
$-6{,}0$ & Prédictions très bien alignées & Excellence \\
\bottomrule
\end{tabular}
\caption{Interprétation de la policy loss}
\end{table}

\subsection{Le signe est l'indicateur de qualité}

\begin{itemize}
    \item \textbf{Loss positive} : l'agent est pénalisé (mauvaises actions renforcées ou bonnes actions ignorées)
    \item \textbf{Loss = 0} : aucun signal d'apprentissage
    \item \textbf{Loss négative} : l'agent est récompensé — ses prédictions sont alignées avec les bonnes actions
    \item \textbf{Loss très négative} : l'agent est fortement en accord avec les bonnes actions
\end{itemize}

Il ne faut \textbf{pas} comparer cette valeur avec les loss positives du supervised learning. En RL, la loss est un \textbf{objectif de maximisation inversé}, pas une erreur à minimiser vers 0.

\section{Question 23 — Loss qui descend mais pas d'amélioration de la reward}

\subsection{Le paradoxe observé}

Dans les pipelines v2 (261011 et 261012), la loss totale descend continuellement, mais la reward moyenne reste plate autour de 0{,}64. Comment est-ce possible ?

\subsection{Explication 1 — La loss et la reward ne mesurent pas la même chose}

La \textbf{loss totale} est un agrégat de 7+ composantes :

\begin{equation}
L = \underbrace{L_{\text{policy}}}_{-5{,}1} + 0{,}3 \cdot \underbrace{L_{\text{contr}}}_{192} + 0{,}5 \cdot \underbrace{L_{\text{value}}}_{0{,}007} + \underbrace{L_{\text{manifold}}}_{?} + \underbrace{L_{\text{prox}}}_{?} - 0{,}1 \cdot \underbrace{H}_{0{,}70} + \ldots
\end{equation}

La loss peut descendre parce que \textbf{certaines composantes s'améliorent} (contrastive loss, value loss, entropie) sans que cela se traduise en \textbf{meilleure sélection de prompts}.

\textbf{Analogie :} un étudiant améliore sa technique de lecture (loss baisse) mais ses notes d'examen ne changent pas (reward stagne) parce que les examens testent des compétences différentes de ce qu'il travaille.

\subsection{Explication 2 — Dominance de la contrastive loss}

La contrastive loss ($\sim$192, pondérée par 0{,}3 $\rightarrow$ contribution de $\sim$57{,}6) domine complètement la loss totale. Quand elle descend un peu (de 193 à 190, soit $-0{,}3 \times 3 = -0{,}9$), la loss totale descend.

Mais cette baisse de la contrastive loss signifie seulement que les embeddings se séparent légèrement mieux — \textbf{pas} que l'agent choisit de meilleurs prompts. La contrastive loss est un objectif auxiliaire de structuration de l'espace, pas directement la performance du système.

\subsection{Explication 3 — La policy loss s'améliore « dans le vide »}

La policy loss passe de $-4{,}2$ à $-5{,}7$. Cela signifie que l'agent aligne mieux ses prédictions avec les actions qu'il \textbf{a déjà prises}. Mais si ces actions étaient déjà sous-optimales (choisies dans un espace mal structuré), mieux les reproduire ne change rien à la qualité.

C'est comme un GPS qui s'améliore pour pointer toujours la même mauvaise direction : il pointe \textbf{plus précisément}, mais toujours au mauvais endroit.

\subsection{Explication 4 — Le plafond de performance}

La reward de $\sim$0{,}64 pourrait être le \textbf{maximum atteignable} avec ce pipeline. Si les embeddings ne séparent pas bien les prompts (contrastive $\sim$192), l'agent ne peut pas faire mieux que le hasard parmi les candidats les plus proches. Il a convergé vers le meilleur choix possible \textbf{étant données les limitations de la représentation}.

\subsection{Explication 5 — Conflit entre composantes}

Les composantes de la loss peuvent être en \textbf{conflit}. Minimiser la contrastive loss pousse les embeddings dans une direction, mais minimiser la policy loss pourrait les pousser dans une direction différente. Le réseau fait un compromis, et ce compromis fait descendre la loss totale sans améliorer la reward.

\subsection{Résumé structuré}

\begin{table}[h]
\centering
\renewcommand{\arraystretch}{1.3}
\begin{tabular}{p{4cm} p{4cm} p{4cm}}
\toprule
\textbf{Ce qui descend} & \textbf{Pourquoi} & \textbf{Impact sur la reward} \\
\midrule
Contrastive loss & Embeddings légèrement mieux structurés & Quasi-nul (trop bruitée) \\
Policy loss & Prédictions mieux alignées & Nul si actions sous-optimales \\
Value loss & Critique plus précis & Indirect seulement \\
Entropie (bonus) & Exploration réduite & Potentiellement négatif \\
\midrule
\textbf{Loss totale} & \textbf{Somme des composantes} & \textbf{Pas d'amélioration} \\
\bottomrule
\end{tabular}
\caption{Déconnexion entre la loss et la reward}
\end{table}

\subsection{Conclusion}

La descente de la loss totale est une \textbf{condition nécessaire mais pas suffisante} pour l'amélioration des performances. Quand la composante dominante (contrastive, $\times$0{,}3 $\times$ 192 = 57{,}6) est « cassée » (valeur anormalement haute), la loss peut descendre en améliorant cette composante sans que le signal informatif ne soit assez fort pour changer le comportement réel de l'agent.

\textbf{Le vrai indicateur de progrès est la reward, pas la loss.}

\end{document}
